\question[4] 
Schreibe eine Funktion, die sich so verhält, wie im folgenden docstring beschrieben.
\begin{lstlisting}
'''
n: int, n > 0
returns: die Summe der echten Teiler von n. Echte Teiler sind alle Teiler der Zahl,
         außer die Zahl selbst. Beispiel: die Summe der echten Teiler
         von 12 ist: 1+2+3+4+6 = 16
 
>>> echteTeilerSumme(12)
16
'''

\end{lstlisting}
\begin{solutionbox}{4cm}
\begin{lstlisting}
def echteTeilerSumme(n):
    summe = 0
    for k in range(1,n):
        if n % k == 0:
            summe +=k
    return summe
\end{lstlisting}
\end{solutionbox}
